% GENERATED FILE. Edit canonical lessons, then run npm run generate:books.
% canonical-source-hash: fnv1a64:2fc77289079fdc9e
% canonical-lessons: LA-C139-cathedra, LA-C139-scamnum, LA-C139-tapete, LA-C139-fax, LA-C139-lanterna

\chapter{Chair with a back, Bench, Carpet, Torch, Lantern}
\label{ch:la-chair-with-a-back-bench-carpet-torch-lantern}
\hlchaptermodality{139}

\begin{chapteropening}
\textbf{By the end of this chapter:} \emph{I can say a chair with a back, a bench, a carpet, a torch and a lantern.}

Complete the last lesson of chapter 139: I can say a chair with a back, a bench, a carpet, a torch and a lantern.
\end{chapteropening}

\section[cathedra, cathedrae]{cathedra, cathedrae — a chair with a back}

\label{lesson:LA-C139-cathedra}



\begin{quote}
Before the new one: say the Latin for a friend, then the Latin for a cupboard.
\end{quote}

\subsection*{You'll want to know: cathedra, cathedrae}
\textbf{cathedra, cathedrae} — \textquotedblleft{}a chair with a back\textquotedblright{}.

\begin{cousinweb}
\textbf{cathedra} was borrowed from Greek \emph{kathedra}, \textquotedblleft{}a seat, a chair\textquotedblright{} (Lewis \& Short), from \emph{kata}, \textquotedblleft{}down\textquotedblright{}, and \emph{hedra}, \textquotedblleft{}seat\textquotedblright{} (EtymDB 2.1). A bishop's \emph{cathedra} gave English \emph{cathedral} its name.
\end{cousinweb}

\begin{grammarlens}[title={What you've built}]
One more word for this chapter.
\end{grammarlens}

\subsection*{Your turn}
\begin{itemize}
\raggedright
  \item \emph{Say it:} \emph{cathedra}
  \item \emph{Say it:} \emph{cathedra}, once more
  \item \emph{Say it:} say \emph{cathedra}
  \item \emph{Recall:} say the Latin for a friend, then the Latin for a cupboard, then say \emph{cathedra} again
\end{itemize}

\begin{tcolorbox}[breakable,colback=teal!4,colframe=teal!35!black,title={Before you move on}]
What does cathedra, cathedrae mean? (\textquotedblleft{}A chair with a back\textquotedblright{}.) Say it once more.
\end{tcolorbox}
\section[scamnum, scamnī]{scamnum, scamnī — a bench}

\label{lesson:LA-C139-scamnum}



\begin{quote}
Before the new one: say the Latin for a cupboard, then the Latin for a chair with a back.
\end{quote}

\subsection*{You'll want to know: scamnum, scamnī}
\textbf{scamnum, scamnī} — \textquotedblleft{}a bench\textquotedblright{}.

\begin{cousinweb}
\textbf{scamnum} is inherited from Proto-Italic \emph{*skabnom}, \textquotedblleft{}bench\textquotedblright{}, usually traced to Proto-Indo-European \emph{*skabh-}, \textquotedblleft{}to prop up\textquotedblright{} (EtymDB 2.1); \emph{scabellum}, \textquotedblleft{}a footstool\textquotedblright{}, is its diminutive (Lewis \& Short).
\end{cousinweb}

\begin{grammarlens}[title={What you've built}]
One more word for this chapter.
\end{grammarlens}

\subsection*{Your turn}
\begin{itemize}
\raggedright
  \item \emph{Say it:} \emph{scamnum}
  \item \emph{Say it:} \emph{scamnum}, once more
  \item \emph{Say it:} say \emph{scamnum}
  \item \emph{Recall:} say the Latin for a cupboard, then the Latin for a chair with a back, then say \emph{scamnum} again
\end{itemize}

\begin{tcolorbox}[breakable,colback=teal!4,colframe=teal!35!black,title={Before you move on}]
What does scamnum, scamnī mean? (\textquotedblleft{}A bench\textquotedblright{}.) Say it once more.
\end{tcolorbox}
\section[tapēte, tapētis]{tapēte, tapētis — a carpet, a rug}

\label{lesson:LA-C139-tapete}



\begin{quote}
Before the new one: say the Latin for a chair with a back, then the Latin for a bench.
\end{quote}

\subsection*{You'll want to know: tapēte, tapētis}
\textbf{tapēte, tapētis} — \textquotedblleft{}a carpet, a rug\textquotedblright{}.

\begin{cousinweb}
\textbf{tapēte} was borrowed from Greek \emph{tapēs}, \emph{tapētos}, \textquotedblleft{}a carpet, a rug\textquotedblright{} (EtymDB 2.1). English \emph{tapestry} belongs to its family, through French.
\end{cousinweb}

\begin{grammarlens}[title={What you've built}]
One more word for this chapter.
\end{grammarlens}

\subsection*{Your turn}
\begin{itemize}
\raggedright
  \item \emph{Say it:} \emph{tapēte}
  \item \emph{Say it:} \emph{tapēte}, once more
  \item \emph{Say it:} say \emph{tapēte}
  \item \emph{Recall:} say the Latin for a chair with a back, then the Latin for a bench, then say \emph{tapēte} again
\end{itemize}

\begin{tcolorbox}[breakable,colback=teal!4,colframe=teal!35!black,title={Before you move on}]
What does tapēte, tapētis mean? (\textquotedblleft{}A carpet, a rug\textquotedblright{}.) Say it once more.
\end{tcolorbox}
\section[fax, facis]{fax, facis — a torch}

\label{lesson:LA-C139-fax}



\begin{quote}
Before the new one: say the Latin for a bench, then the Latin for a carpet.
\end{quote}

\subsection*{You'll want to know: fax, facis}
\textbf{fax, facis} — \textquotedblleft{}a torch\textquotedblright{}.

\begin{cousinweb}
\textbf{fax} has no sure origin; EtymDB 2.1 records a proposed Proto-Indo-European \emph{*gwheh2k-}, \textquotedblleft{}to shine\textquotedblright{}, but the link is not settled.
\end{cousinweb}

\begin{grammarlens}[title={What you've built}]
One more word for this chapter.
\end{grammarlens}

\subsection*{Your turn}
\begin{itemize}
\raggedright
  \item \emph{Say it:} \emph{fax}
  \item \emph{Say it:} \emph{fax}, once more
  \item \emph{Say it:} say \emph{fax}
  \item \emph{Recall:} say the Latin for a bench, then the Latin for a carpet, then say \emph{fax} again
\end{itemize}

\begin{tcolorbox}[breakable,colback=teal!4,colframe=teal!35!black,title={Before you move on}]
What does fax, facis mean? (\textquotedblleft{}A torch\textquotedblright{}.) Say it once more.
\end{tcolorbox}
\section[lanterna, lanternae]{lanterna, lanternae — a lantern}

\label{lesson:LA-C139-lanterna}



\begin{quote}
Before the new one: say the Latin for a carpet, then the Latin for a torch.

Count on past the ninth hour: \emph{nōnus}, \emph{decimus}. Say \emph{decima hōra}, then say whether it is \emph{ante} or \emph{post merīdiem}.
\end{quote}

\subsection*{You'll want to know: lanterna, lanternae}
\textbf{lanterna, lanternae} — \textquotedblleft{}a lantern\textquotedblright{}.

\begin{cousinweb}
\textbf{lanterna} comes from Greek \emph{lamptēr}, \textquotedblleft{}a torch, a lamp\textquotedblright{} (Lewis \& Short), from \emph{lampō}, \textquotedblleft{}to shine\textquotedblright{}, reshaped with the ending of \emph{lucerna}, \textquotedblleft{}lamp\textquotedblright{}.
\end{cousinweb}

\begin{grammarlens}[title={What you've built}]
One more word for this chapter.
\end{grammarlens}

\subsection*{Your turn}
\begin{itemize}
\raggedright
  \item \emph{Say it:} \emph{lanterna}
  \item \emph{Say it:} \emph{lanterna}, once more
  \item \emph{Say it:} say \emph{lanterna}
  \item \emph{Recall:} say the Latin for a carpet, then the Latin for a torch, then say \emph{lanterna} again
\end{itemize}

\begin{tcolorbox}[breakable,colback=teal!4,colframe=teal!35!black,title={Before you move on}]
What does lanterna, lanternae mean? (\textquotedblleft{}A lantern\textquotedblright{}.) Say it once more.
\end{tcolorbox}
